\documentclass[11pt]{article}
\usepackage{mathblatt}
% Präambel des Setzers (werkzeuge/setzer-praeambel.tex, wortgleich); dazu \gteil (Muster T2U)
\makeatletter
\setlength{\mbpfbe}{0mm}% keine Punkte (pruefung.md § 5: nur beim Prüfstein)
\newlength{\gEin}\setlength{\gEin}{0pt}
\let\gSchrift\relax
\renewcommand{\pfkreuzzeile}[1]{\par\noindent\hangindent3.2em\hangafter1\hspace*{1.6em}$\square$\ #1\par}
\newcommand{\gkreuz}[1]{$\square$\ #1\hspace{2em}}
\newcommand{\gkreuzzeile}[2]{\par\noindent\hangindent2.2em\hangafter1\mbox{#1\,$\square$}\ #2\par}% Buchstabe, Option: eine je Zeile (neben Skizze)
% links, was man ansieht, rechts, was man tut; Spaltengrenze überall an derselben Stelle (gemeinsam.md „Satz“)
\newsavebox{\gbox}\newlength{\gLinks}
% v0.10 (Lehrer 13.10.): keine starre Halbierung – die linke Spalte ist so breit wie die Figur (+ 6 mm), die
% Antwort steht direkt daneben; ist die Figur breiter als 0,62 Zeile, steht die Antwort darunter
\newcommand{\gzwei}[2]{\par\smallskip\noindent\sbox{\gbox}{#1}%
  \setlength{\gLinks}{\dimexpr\wd\gbox+6mm\relax}%
  \ifdim\gLinks>0.62\linewidth
    \usebox{\gbox}\par\smallskip\noindent\begin{minipage}{\linewidth}#2\end{minipage}%
  \else
    \begin{minipage}[t]{\gLinks}\vspace{0pt}\usebox{\gbox}\end{minipage}%
    \begin{minipage}[t]{\dimexpr\linewidth-\gLinks\relax}\vspace{0pt}\raggedright #2\end{minipage}%
  \fi\par}
% v0.9 (12.10., pruefung.md § 5 „Marke mit B, Bedingungen rechts“): drittes Argument von pfaufg = Bedingungen
% der Prüfung (v0.10: \pfohnetr = TR durchgestrichen, „FOR“), klein grau im rechten Rand auf der Grundlinie der Nummer, Zeilen mit \\
\newcommand{\gbedrand}[2]{\ifblank{#2}{}{\rlap{\hspace*{\dimexpr#1+2mm\relax}%
  \parbox[t]{15mm}{\raggedright\scriptsize\leavevmode\color{mbgrau}\strut#2}}}}
% Teil eines Bündels (a, b …): Marke links im Rand, Bedingungen rechts im Rand, auf der Zeile des Teils
\newcommand{\gteilrand}[2]{\llap{\mbpfrandmarke{#1}\hspace*{\mbpfnr}}\gbedrand{\linewidth}{#2}}
% Teil eines Bündels: Buchstabe in eigener Spalte, Inhalt eingerückt (\gEin = Einzug, damit die Spaltengrenze von
% \gzwei an derselben Stelle bleibt)
\newenvironment{gbteil}[3]{\par\smallskip\noindent\gteilrand{#2}{#3}\makebox[\mbpfteil][l]{\textbf{#1}}%
  \begin{minipage}[t]{\dimexpr\linewidth-\mbpfteil\relax}\setlength{\gEin}{\mbpfteil}\raggedright\strut\ignorespaces}%
 {\end{minipage}\par}
% Grafik höchstens so breit wie die Spalte (Bündel nebeneinander, Koordinatensystem ohne Karo)
\newsavebox{\gpbox}
\newcommand{\gpassend}[1]{\sbox{\gpbox}{#1}\ifdim\wd\gpbox>\linewidth\typeout{GPASSEND: verkleinert \the\wd\gpbox\space auf \the\linewidth}\resizebox{\linewidth}{!}{\usebox{\gpbox}}\else\usebox{\gpbox}\fi}
% Skizzenfeld (v0.10, Lehrer 13.10.): Karo ab dem linken Rand über die Zeile, „Skizze“ klein grau oben links im Feld
% (wie \pfskizzenfeld der Vorlage); der rechte Rand bleibt den Bedingungen
\newcommand{\gskizze}[1]{\par\smallskip\noindent\begin{tikzpicture}
  \clip (0,0) rectangle ({\linewidth-1mm},{#1*10mm});
  \draw[black!16,line width=0.3pt,step=5mm] (0,0) grid ({\linewidth},{#1*10mm});
  \node[anchor=north west,inner sep=2pt,font=\scriptsize,text=mbgrau,fill=white] at (0,{#1*10mm}) {Skizze};
  \end{tikzpicture}\par}
\RenewDocumentEnvironment{pfaufg}{m m m}{%
  \begin{lrbox}{\mbaufgabebox}\begin{minipage}[t]{\linewidth}%
  \gSchrift\noindent\hspace*{\gEin}\mbpfrandmarke{#1}\gbedrand{\linewidth-\gEin}{#3}%
  \begin{minipage}[t]{\mbpfnr}\strut\textbf{#2}\end{minipage}%
  \begin{minipage}[t]{\dimexpr\linewidth-\gEin-\mbpfnr\relax}\raggedright\strut\ignorespaces}%
 {\end{minipage}%
  \end{minipage}\end{lrbox}%
  \setlength{\mbaufgabehoehe}{\dimexpr\ht\mbaufgabebox+\dp\mbaufgabebox\relax}%
  \ifdefempty{\gkopfnach}{\par\addvspace{5pt}\Needspace*{\mbaufgabehoehe}}{%
    % Stufenkopf, der mit dieser (ersten) Aufgabe zusammenbleibt: Platz für beide verlangen, dann den Kopf setzen
    % (12.10., Quadratische: ein Kopf stand allein unten, die Aufgabe mit Koordinatensystem auf der nächsten Seite)
    \par\Needspace*{\dimexpr\mbaufgabehoehe+40pt\relax}\gkopfnach\gdef\gkopfnach{}\par\addvspace{5pt}}%
  \noindent\usebox{\mbaufgabebox}\par\addvspace{5pt}}
% Stufenkopf vormerken: \gkopfvor{<Kopf>} – die nächste pfaufg setzt ihn, sobald beide zusammen Platz haben
\def\gkopfnach{}
\newcommand{\gkopfvor}[1]{\gdef\gkopfnach{#1}}
% Teilaufgabe mit Skizze links in gzwei: Buchstabe in eigener Spalte, darunter Text und Skizze
\newcommand{\gteil}[2]{\makebox[\mbpfteil][l]{\textbf{#1}}\parbox[t]{\dimexpr0.5\textwidth-\gEin-\mbpfnr-5mm-\mbpfteil\relax}{\raggedright #2}}
% Stufenkopf der Vorlage (mathblatt.sty 2026-10-12: Linie, Abstand, 13 pt, \pfjahre) ohne ihren Vorhalt von
% 14 Zeilen: der Setzer fordert vor dem Kopf 0,16 Seitenhöhe (Kopf + erste Aufgabe), mehr nicht
\setlength{\mbpfstufevorhalt}{0pt}
\makeatother
\newcommand{\karo}[1]{\par\smallskip\noindent\begin{tikzpicture}
  \clip (0,0) rectangle ({\linewidth-1mm},{#1*10mm});
  \draw[black!16,line width=0.3pt,step=5mm] (0,0) grid ({\linewidth},{#1*10mm});
  \end{tikzpicture}\par}
\newcommand{\kk}[1]{$\square$\,#1\hspace{0.9em}}
\newcommand{\linie}[1]{\,{\color{black!45}\rule[-1pt]{#1}{0.4pt}}\,}
\newcommand{\pkt}{\textperiodcentered{}}
% v0.14 (14.10.): Stufenkopf mit fünf Punkten (Option = Jahre für die Bestellkarte) oder „zuletzt 19“
\NewDocumentCommand{\gpunkte}{o m}{#2}
\newcommand{\gzuletzt}[1]{zuletzt #1}
% v0.7 (12.10.): Schülerfrage unter dem Stufenkopf; Ankreuzoption als Bild (Buchstabe und Kästchen darunter)
\newcommand{\gfrage}[1]{\par\nobreak\noindent\hspace*{\mbpfnr}{\small\itshape #1}\par\nobreak\addvspace{2pt}}
\newlength{\gbildbreite}\setlength{\gbildbreite}{0.32\linewidth}% Setzer setzt sie je Reihe (vier Bilder: 0,24)
\newcommand{\gbildoption}[3]{\begin{minipage}[t][#3][b]{\gbildbreite}\centering #2\par\smallskip\mbox{#1\,$\square$}\end{minipage}}
% Lineare (12.10.): Bild mit Buchstabe ohne Kästchen (Diagramm zuordnen); Aussage mit „wahr / falsch“ zum Ankreuzen
\newcommand{\gbildname}[3]{\begin{minipage}[t][#3][b]{\gbildbreite}\centering #2\par\smallskip\mbox{\textbf{#1}}\end{minipage}}
\newcommand{\gwahrfalsch}[2]{\par\noindent\hangindent2.2em\hangafter1\makebox[2.2em][l]{(#1)}#2\hfill\mbox{wahr $\square$\quad falsch $\square$}\par}
\makeatletter\setlength{\mbpfbe}{0mm}\makeatother
\begin{document}

\pfheftstil
\pflaufkopf[P10 \pkt{} mehr]{Fläche und Umfang}{PG40-7}\pfseitevon
\pfheftkopf{Fläche und Umfang}{mehr}
\pfstufenzeile{Formeln umstellen \pkt{} Formeln \pkt{} Fläche gesucht}
% 2020-OS-B1d-m1 (Schwere I, rang 12)
\begin{pfaufg}{nach~P10\\’20~\pkt{}~1d}{1.}{}
Ein Quadrat hat den Flächeninhalt 81 cm². Gib seine Seitenlänge an.
\karo{2}
\fusshilfe{\mbox{1:~9 cm}}
\end{pfaufg}
% 2019-OS-B1e-m1 (Schwere I, rang 13)
\begin{pfaufg}{nach~P10\\’19~\pkt{}~1e}{2.}{}
Gib eine Gleichung für den Flächeninhalt A der gesamten Fläche an.
\par\smallskip\noindent\gpassend{\begin{tikzpicture}[line width=0.8pt,font=\small]\fill[black!20] (0,0) rectangle (1.2,1.4);\draw[black!45,line width=0.3pt,step=0.2] (1.2,0) grid (3.0,1.4);\draw (0,0) rectangle (3.0,1.4);\draw (1.2,0)--(1.2,1.4);\node[below] at (0.6,0) {$a$};\node[below] at (2.1,0) {$b$};\node[right] at (3.0,0.7) {$c$};\end{tikzpicture}}\par
\hfill A $=$\linie{30mm}
\fusshilfe{\mbox{2:~(a + b) · c}}
\end{pfaufg}
% 2018-OS-B1f-m1 (Schwere I, rang 14)
\begin{pfaufg}{nach~P10\\’18~\pkt{}~1f}{3.}{}
Ein Rechteck hat den Umfang u = 30 cm; die Seite a ist 11 cm lang. Gib die Länge der Seite b an.
\par\smallskip\noindent\gpassend{\begin{tikzpicture}[line width=0.8pt,font=\small]\draw (0,0) rectangle (3.3,1.2);\node[below] at (1.65,0) {$a = 11$ cm};\node[right] at (3.3,0.6) {$b$};\end{tikzpicture}}\par
\karo{2}
\fusshilfe{\mbox{3:~4 cm}}
\end{pfaufg}
% 2017-OS-K3a-m1 (Schwere I, rang 15)
\begin{pfaufg}{nach~P10\\’17~\pkt{}~3a}{4.}{}
Familie Sommer hat einen Swimmingpool gebaut. Seine Grundfläche ist ein Rechteck, an dessen kurze Seiten je ein Halbkreis angesetzt ist. Die kurze Seite des Rechtecks ist a = 6 m lang, die lange Seite b = 11 m. Gib an, aus welchen Teilflächen die Grundfläche des Pools besteht.
\par\smallskip\noindent\gpassend{\begin{tikzpicture}[line width=0.8pt,font=\small]\filldraw[fill=black!12] (0,0)--(2.2,0) arc (-90:90:0.6)--(0,1.2) arc (90:270:0.6)--cycle;\draw[dashed,line width=0.5pt] (0,0)--(0,1.2) (2.2,0)--(2.2,1.2);\node[left] at (-0.6,0.6) {$a = 6$ m};\node[below] at (1.1,0) {$b = 11$ m};\end{tikzpicture}}\par
\hfill Ein Rechteck und zwei Halbkreise (zusammen ein Kreis mit Durchmesser a) $=$\linie{30mm}
\fusshilfe{\mbox{4:~Ein Rechteck und zwei Halbkreise (zusammen ein Kreis mit Durchmesser a)}}
\end{pfaufg}
% 2016-OS-K3b-m1 (Schwere I, rang 16)
\begin{pfaufg}{nach~P10\\’16~\pkt{}~3b}{5.}{}
Beim Kugelstoßen steht man in einem Abwurfring; das ist ein Kreis mit 2,30 m Durchmesser. Ermittle den Flächeninhalt dieses Kreises.
\karo{2}
\fusshilfe{\mbox{5:~4,15 m²}}
\end{pfaufg}
% 2014-OS-B1g-m1 (Schwere I, rang 17)
\begin{pfaufg}{nach~P10\\’14~\pkt{}~1g}{6.}{}
Die graue Figur ist aus gleich großen Quadraten mit der Seitenlänge a zusammengesetzt. Gib einen Term für ihren Umfang u an.
\par\smallskip\noindent\gpassend{\begin{tikzpicture}[scale=0.4,line width=0.8pt,font=\small]\draw[black!20,line width=0.3pt] (-1,-1) grid (8,6);\fill[black!30] (0,2) rectangle (7,3);\fill[black!30] (3,0) rectangle (4,5);\draw (0,2)--(3,2)--(3,0)--(4,0)--(4,2)--(7,2)--(7,3)--(4,3)--(4,5)--(3,5)--(3,3)--(0,3)--cycle;\node[left] at (0,2.5) {$a$};\end{tikzpicture}}\par
\karo{2}
\fusshilfe{\mbox{6:~24 · a}}
\end{pfaufg}
% 2020-OS-K7b-m1 (Schwere II, rang 18)
\begin{pfaufg}{nach~P10\\’20~\pkt{}~7b}{7.}{}
Im Dreieck ABC ist bei C ein rechter Winkel; BC ist 15 m und AC 28 m lang. Der Punkt Q liegt auf AC. Das Dreieck QBC hat den Flächeninhalt 120 m². Berechne die Länge der Strecke QC.
\par\smallskip\noindent\gpassend{\begin{tikzpicture}[line width=0.8pt,font=\small]\coordinate (C) at (0,2.5);\coordinate (B) at (1.35,2.5);\coordinate (A) at (0,0);\coordinate (Q) at (0,1.07);\draw (A)--(B)--(C)--cycle;\draw[dashed,line width=0.5pt] (Q)--(B);\pfrw[0.25]{(C)}{270}\node[above left] at (C) {$C$};\node[above right] at (B) {$B$};\node[below left] at (A) {$A$};\node[left] at (Q) {$Q$};\end{tikzpicture}}\par
\karo{3}
\fusshilfe{\mbox{7:~16 m}}
\end{pfaufg}
% 2019-OS-K4c-m1 (Schwere II, rang 19)
\begin{pfaufg}{nach~P10\\’19~\pkt{}~4c}{8.}{}
Auf einem Mast sitzt ein Prisma, dessen Grund- und Deckfläche gleichseitige Dreiecke mit 1,80 m Seitenlänge sind; die Höhe dieser Dreiecke beträgt 1,56 m. Das Prisma ist 1,50 m hoch. Ermittle den Flächeninhalt der Deckfläche. Durch ein Loch ist Regenwasser in das Prisma gelaufen. Zeige mit einer Rechnung, dass etwa 2,1 m³ Wasser hineinpassen. Leer wiegt das Prisma 250 kg; 1 m³ Regenwasser wiegt 1 000 kg. Berechne, wie schwer das Prisma ist, wenn ein Zehntel seines Innenraums (10 \%) voll Wasser ist.
\par\smallskip\noindent\gpassend{\begin{tikzpicture}[line width=0.8pt,font=\small]\filldraw[fill=black!15] (0,0)--(2.2,0)--(1.1,1.905)--cycle;\draw[dashed,line width=0.5pt] (1.1,0)--(1.1,1.905);\pfrw[0.2]{(1.1,0)}{0}\node[below] at (1.1,0) {$1{,}80$ m};\node[left,inner sep=1.5pt] at (1.1,0.75) {$h$};\end{tikzpicture}}\par
\karo{4}
\fusshilfe{\mbox{8:~1,40 m²; 2,1 m³; 461 kg}}
\end{pfaufg}
% 2018-OS-K6a-m1 (Schwere II, rang 20)
\begin{pfaufg}{nach~P10\\’18~\pkt{}~6a}{9.}{}
Ein runder Turm mit 5,8 m Durchmesser steht auf einer Rasenfläche, die 18,0 m × 11,0 m groß ist. Der Rasen um den Turm herum soll neu angelegt werden. Berechne die Fläche des neuen Rasens.
\par\smallskip\noindent\gpassend{\begin{tikzpicture}[line width=0.8pt]\draw[dotted] (0,0) rectangle (3.6,2.2);\filldraw[fill=black!15] (1.8,1.1) circle (0.58);\end{tikzpicture}}\par
\karo{2}
\fusshilfe{\mbox{9:~171,6 m²}}
\end{pfaufg}
% 2017-OS-K3b-m1 (Schwere II, rang 21)
\begin{pfaufg}{nach~P10\\’17~\pkt{}~3b}{10.}{}
Familie Sommer hat einen Swimmingpool gebaut. Seine Grundfläche ist ein Rechteck, an dessen kurze Seiten je ein Halbkreis angesetzt ist. Die kurze Seite des Rechtecks ist a = 6 m lang, die lange Seite b = 11 m. Berechne den Flächeninhalt der Grundfläche des Pools.
\par\smallskip\noindent\gpassend{\begin{tikzpicture}[line width=0.8pt,font=\small]\filldraw[fill=black!12] (0,0)--(2.2,0) arc (-90:90:0.6)--(0,1.2) arc (90:270:0.6)--cycle;\draw[dashed,line width=0.5pt] (0,0)--(0,1.2) (2.2,0)--(2.2,1.2);\node[left] at (-0.6,0.6) {$a = 6$ m};\node[below] at (1.1,0) {$b = 11$ m};\end{tikzpicture}}\par
\karo{2}
\fusshilfe{\mbox{10:~94,3 m²}}
\end{pfaufg}
% 2015-OS-K6c-m1 (Schwere II, rang 22)
\begin{pfaufg}{nach~P10\\’15~\pkt{}~6c}{11.}{}
Mia baut das Modell einer Pyramide mit quadratischer Grundfläche: Die Grundkante ist 0,90 m lang, die Pyramide ist 0,70 m hoch. Die Seitenflächen sägt sie aus dünnem Sperrholz. Weise nach, dass sie für eine Seitenfläche etwa 0,37 m² Sperrholz braucht.
\karo{3}
\fusshilfe{\mbox{11:~0,83 m; 0,37 m²}}
\end{pfaufg}
% 2014-OS-K5c-m1 (Schwere II, rang 23)
\begin{pfaufg}{nach~P10\\’14~\pkt{}~5c}{12.}{}
Eine Kühlvitrine hat die Form eines Prismas. Ihre Grundfläche ist ein gleichschenkliges Trapez mit den parallelen Seiten 108 cm und 84 cm und dem Flächeninhalt 3 360 cm². Die Höhe h dieses Trapezes ist die Tiefe der Vitrine. Weise nach, dass Kuchenplatten mit 30 cm Durchmesser in die Vitrine passen.
\par\smallskip\noindent\gpassend{\begin{tikzpicture}[line width=0.8pt,font=\small]\draw (0,1.05)--(3.24,1.05)--(2.88,0)--(0.36,0)--cycle;\draw[dashed,line width=0.5pt] (1.62,0)--(1.62,1.05);\pfrw[0.18]{(1.62,0)}{0}\pfrw[0.18]{(1.62,1.05)}{270}\node[left] at (1.62,0.52) {$h$};\node[above] at (1.62,1.05) {$108$ cm};\node[below] at (1.62,0) {$84$ cm};\node[right] at (3.06,0.52) {$37$ cm};\end{tikzpicture}}\par
\karo{2}
\fusshilfe{\mbox{12:~35 cm}}
\end{pfaufg}
% 2014-OS-K5d-m1 (Schwere II, rang 24)
\begin{pfaufg}{nach~P10\\’14~\pkt{}~5d}{13.}{}
Die Rückwand derselben Vitrine ist ein Rechteck, 108 cm breit und 150 cm hoch. Sie soll aus einer rechteckigen Platte gesägt werden, die 1,25 m breit ist und einen Flächeninhalt von 1,8 m² hat. Entscheide mit einer Rechnung, ob die Platte dafür reicht.
\karo{2}
\fusshilfe{\mbox{13:~1,44 m hoch}}
\end{pfaufg}
% 2015-OS-K5d-m1 (Schwere III, rang 25)
\begin{pfaufg}{nach~P10\\’15~\pkt{}~5d}{14.}{}
Im Dreieck AED ist bei E ein rechter Winkel. Die Hypotenuse AD ist 3,9 cm lang, der Winkel bei D beträgt $\delta$ = 28°. Der Flächeninhalt des Dreiecks AED soll berechnet werden. Gib an, welche Seiten und Winkel du dazu brauchst. Schreibe die Rechenschritte in einer sinnvollen Reihenfolge auf.
\par\smallskip\noindent\gpassend{\begin{tikzpicture}[line width=0.8pt,font=\small]\coordinate (E) at (0,0);\coordinate (D) at (3.8,0);\coordinate (A) at (0,2.02);\draw (E)--(D)--(A)--cycle;\pfrw[0.3]{(E)}{0}\pfwinkel[1.5][0.72]{(D)}{152}{180}{28^\circ}\node[below left] at (E) {$E$};\node[below right] at (D) {$D$};\node[above left] at (A) {$A$};\node[above right] at (1.9,1.01) {$3{,}9$ cm};\end{tikzpicture}}\par
\karo{5}
\fusshilfe{\mbox{14:~28°; 1,83 cm; 3,44 cm; 3,15 cm²}}
\end{pfaufg}
\par\vfill\noindent{\scriptsize\color{mbgrau}Originale: P10 2020 \pkt{} 1d, P10 2019 \pkt{} 1e, P10 2018 \pkt{} 1f, P10 2017 \pkt{} 3a, P10 2016 \pkt{} 3b, P10 2014 \pkt{} 1g, P10 2020 \pkt{} 7b, P10 2019 \pkt{} 4c, P10 2018 \pkt{} 6a, P10 2017 \pkt{} 3b, P10 2015 \pkt{} 6c, P10 2014 \pkt{} 5c, P10 2014 \pkt{} 5d, P10 2015 \pkt{} 5d\par}
\end{document}